\section*{Bab \ref{ch:b3:fouriertransform} --- Transformasi
Fourier}

\begin{solution}{exo:b3:fouriertransform:1}
$\widehat{\mathbf 1_{\intcc{-a}a}}(\xi) =
\int_{-a}^a\eu^{-\iu\xi x}\dd x = \frac{2\sin(a\xi)}{\xi}$
(dengan nilai $2a$ di $0$).
$\widehat{\eu^{-a\abs x}}(\xi) = \int_0^\infty\eu^{-(a +
\iu\xi)x} + \eu^{-(a - \iu\xi)x}\,\dd x = \frac1{a + \iu\xi} +
\frac1{a - \iu\xi} = \frac{2a}{a^2 + \xi^2}$.
Untuk tendanya: $\max(0, 1 - \abs x) = \mathbf 1_{\intcc{-1/2}{1/2}} *
\mathbf 1_{\intcc{-1/2}{1/2}}$, sehingga transformasinya
$\bigl(\frac{2\sin(\xi/2)}\xi\bigr)^2 =
\bigl(\frac{\sin(\xi/2)}{\xi/2}\bigr)^2$.
Yang terakhir: $\frac{2a}{a^2+\xi^2} \in L^1$, sehingga pembalikan
(\cref{thm:b3:fouriertransform:inversion}(2)) yang diterapkan pada
$\eu^{-a\abs x}$ memberikan, setelah variabelnya dinamai ulang,
\[
\widehat{\Bigl(\frac1{x^2 + a^2}\Bigr)}(\xi) =
\frac{\pi}{a}\,\eu^{-a\abs\xi} .
\]
\end{solution}

\begin{solution}{exo:b3:fouriertransform:2}
Dari \cref{prop:b3:fouriertransform:basic}:
$\widehat{f(\cdot - a)} = \eu^{-\iu a\xi}\hat f(\xi)$;
$\widehat{f\cos(b\cdot)} = \frac12\bigl(\hat f(\xi - b) + \hat
f(\xi + b)\bigr)$;
$\widehat{f(a\cdot + b)}(\xi) = \frac1a\,\eu^{\iu
b\xi/a}\,\hat f(\xi/a)$ (dengan $a > 0$);
$\widehat{\overline{f(-\cdot)}} = \overline{\hat f}$;
$\widehat{f * f} = \hat f^2$. Pada fungsi Gauss
(sebab $\widehat{\eu^{-x^2}} = \sqrt\pi\eu^{-\xi^2/4}$) setiap aturannya berupa
pemeriksaan satu baris --- misalnya $\eu^{-(x-a)^2}$ bertransformasi
$\sqrt\pi\,\eu^{-\iu a\xi}\eu^{-\xi^2/4}$, yang dibenarkan
perhitungan langsungnya (lengkapkan kuadratnya).
\end{solution}

\begin{solution}{exo:b3:fouriertransform:3}
(a) Fungsi $h = \mathbf 1_{\intcc{-1}1}*\mathbf 1_{\intcc{-1}1}$ mempunyai
$\hat h = \bigl(\frac{2\sin\xi}\xi\bigr)^2 \in L^1$; lalu pembalikan
di $x = 0$, tempat $h(0) = \lambda(\intcc{-1}1\cap\intcc{-1}1)
= 2$:
\[
2 = \frac1{2\pi}\int_\R\Bigl(\frac{2\sin\xi}\xi\Bigr)^2
\dd\xi
\ \Longrightarrow\
\int_\R\Bigl(\frac{\sin\xi}\xi\Bigr)^2\dd\xi = \pi ,
\]
yang konsisten dengan $\int_0^\infty\frac{\sin^2}{\xi^2} =
\frac\pi2$ (\cref{pb:b3:lebesgue:1}).

(b) Plancherel untuk $f = \eu^{-\abs x}$: $\int\abs{\hat f}^2 =
2\pi\int\abs f^2$ berbunyi $\int\frac{4\,\dd\xi}{(1 + \xi^2)^2} =
2\pi\int\eu^{-2\abs x}\dd x = 2\pi$: sehingga
$\int_\R\frac{\dd\xi}{(1+\xi^2)^2} = \frac\pi2$.
\end{solution}

\begin{solution}{exo:b3:fouriertransform:4}
(a) $g_t(x) = \frac1{2\sqrt{\pi t}}\eu^{-x^2/4t}$: menurut
\cref{ex:b3:fouriertransform:gaussian} dengan $a = \frac1{4t}$,
$\hat g_t(\xi) = \frac1{2\sqrt{\pi t}}\sqrt{4\pi
t}\,\eu^{-t\xi^2} = \eu^{-t\xi^2}$.
(b) $\widehat{g_t * g_s} = \eu^{-t\xi^2}\eu^{-s\xi^2} =
\widehat{g_{t+s}}$, dan transformasinya injektif pada $L^1$
(\cref{thm:b3:fouriertransform:inversion}(3)): sehingga $g_t * g_s =
g_{t+s}$.
(c) $\norm{g_t}_1 = 1$ (lewat integral Gaussnya); dan
$\norm{g_t}_2^2 = \frac1{4\pi t}\int\eu^{-x^2/2t}\dd x =
\frac{\sqrt{2\pi t}}{4\pi t} = \frac1{\sqrt{8\pi t}}$.
\end{solution}

\begin{solution}{exo:b3:fouriertransform:5}
Untuk $f$ yang real dan genap: $\hat f(\xi) = \int f\cos(\xi x)\dd x$ (sebab
bagian sinusnya meniadakan diri): jadi real dan genap. Untuk yang gasal: $\hat f(\xi) =
-\iu\int f\sin(\xi x)$: jadi gasal dan murni imajiner. Sedangkan $\hat f(0) =
\int f$: yakni massa totalnya. Jika $f \geq 0$: $\abs{\hat f(\xi)}
\leq \int\abs f = \int f = \hat f(0)$, sehingga supremumnya tercapai
di $0$. Untuk sebuah \omterm{ex:b3:lebesgue:gamma}{kepadatan} peluang, $\hat f(-\xi)$ adalah
fungsi karakteristik \cref{ch:b3:clt}: yang bermodulus $\leq 1$
di mana-mana dan $= 1$ di titik asalnya.
\end{solution}

\begin{solution}{exo:b3:fouriertransform:6}
(i) Keinjektifannya adalah
\cref{thm:b3:fouriertransform:inversion}(3); \omterm{def:b3:topology:continuity}{kekontinuannya} adalah
$\norm{\hat f}_\infty \leq \norm f_1$ (dengan nilainya di
$\mathcal C_0$ menurut Riemann--Lebesgue); dan $\mathcal C_0$ tertutup
dalam norma supnya (sebab limit seragam fungsi yang lenyap di tak hingga
lenyap di tak hingga): jadi Banach.

(ii) Bijeksi \omterm{def:b3:topology:continuity}{kontinu} antara ruang Banach berbalikan
\omterm{def:b3:topology:continuity}{kontinu} (\cref{thm:b3:banach:openmapping}): sehingga
akan ada $C$ dengan $\norm f_1 \leq C\norm{\hat f}_\infty$.

(iii) Misalkan $f_n(x) = \frac{\sin x}x\cdot\frac{\sin(x/n)}{x/n}$:
yakni hasil kali dua fungsi $L^2$, yang bernilai $O(x^{-2})$ di
tak hingga, sehingga $f_n \in L^1\cap L^2$. Karena
$\bigl(\frac{\sin(ax)}{ax}\bigr)$ bertransformasi-$L^2$
$\frac\pi a\mathbf 1_{\intcc{-a}a}$, rumus hasil kalinya
$\widehat{gh} = \frac1{2\pi}\hat g * \hat h$ (yang sahih untuk $g, h
\in L^2$ dengan $gh \in L^1$; periksalah pada fungsi Schwartz lewat
Fubini lalu perluas lewat kekontinuan-$L^2$ kedua ruasnya lewat
Plancherel) memberikan
\[
\hat f_n = \frac1{2\pi}\,\bigl(\pi\mathbf
1_{\intcc{-1}1}\bigr) * \bigl(\pi n\,\mathbf
1_{\intcc{-1/n}{1/n}}\bigr):
\]
yakni sebuah trapesium setinggi $\frac{\pi n}2\cdot\frac2n = \pi$:
sehingga $\norm{\hat f_n}_\infty = \pi$ untuk setiap $n$. Padahal pada $[1,
n]$, $\frac{\sin(x/n)}{x/n} \geq \sin 1 > 0$, sehingga
\[
\norm{f_n}_1 \geq \sin 1\int_1^n\frac{\abs{\sin x}}x\dd x
\geq c\ln n
\]
(lewat pencacahan lengkungannya, seperti pada \cref{thm:b3:banach:fourierdiverge}).
Jadi batas $\norm{f_n}_1 \leq C\pi$ gagal untuk $n$ yang besar: sehingga tak
surjektif. (Petanya merupakan subruang $\mathcal C_0$ yang padat --- lewat
argumen bertipe Stone--Weierstrass --- tetapi sejati.)
\end{solution}

\begin{solution}{exo:b3:fouriertransform:7}
Jika $\hat f(\xi) = O(\abs\xi^{-k-1-\delta})$: maka $\xi^j\hat f
\in L^1$ untuk $0 \leq j \leq k$ (yang \omterm{def:b3:lebesgue:l1}{terintegralkan} di tak hingga menurut
peluruhannya, dan secara lokal menurut \omterm{def:b3:topology:continuity}{kekontinuan} $\hat f$). Lalu pembalikan
(\cref{thm:b3:fouriertransform:inversion}(2)) menyajikan $f$
hampir di mana-mana lewat $x \mapsto \frac1{2\pi}\int\hat f(\xi)\eu^{\iu
x\xi}\dd\xi$, dan pendiferensialan di bawah integralnya
(dengan pendominasi $\abs{\xi^j\hat f}$) membuat wakil itu
$\mathcal C^k$. Sebaliknya untuk $f \in \mathcal C_c^k$:
dengan mengiterasikan \cref{prop:b3:fouriertransform:basic}(4),
$(\iu\xi)^k\hat f = \widehat{f^{(k)}}$, sehingga $\abs{\hat f} \leq
\norm{f^{(k)}}_1\abs\xi^{-k}$. Untuk fungsi tendanya: ia \omterm{def:b3:topology:continuity}{kontinu} dengan
pendukung \omterm{def:b3:topology:compact}{kompak} (sehingga $k = 0$: transformasinya terbatas), dan
transformasinya $\sim \xi^{-2} = O(\abs\xi^{-0-1-1})$ mengembalikan, lewat
arah pertamanya, sebuah wakil $\mathcal C^0$ --- keduanya
tajam: sebab tendanya bukan $\mathcal C^1$, dan transformasinya
meluruh tak lebih cepat daripada $\xi^{-2}$.
\end{solution}

\begin{solution}{exo:b3:fouriertransform:8}
Lewat pengintegralan parsial (sebab $f \in \mathcal S$; sehingga suku batasnya
lenyap):
\[
1 = \int f^2 = \bigl[xf^2\bigr]_{-\infty}^{\infty} - \int
x\,(f^2)' = -2\int xff' \leq 2\,\norm{xf}_2\,\norm{f'}_2 .
\]
Lalu Plancherel dan $\widehat{f'} = \iu\xi\hat f$:
$\norm{f'}_2^2 = \frac1{2\pi}\int\xi^2\abs{\hat f}^2$. Dengan mengkuadratkan
tampilannya:
\[
\frac14 \leq \norm{xf}_2^2\cdot\frac1{2\pi}
\int\xi^2\abs{\hat f}^2\dd\xi .
\]
Kesamaannya menuntut kesamaan pada Cauchy--Schwarz: yakni $f' =
\lambda xf$ dengan $\lambda < 0$ (menurut keterintegralannya), yakni $f(x) =
c\,\eu^{\lambda x^2/2}$: jadi fungsi Gauss. Sebuah sinyal yang terpusat pada
$x$ (dengan $\norm{xf}_2$ yang kecil) haruslah berspektrum tersebar,
dan sebaliknya: yakni asas ketakpastiannya.
\end{solution}

\begin{solution}{exo:b3:fouriertransform:9}
Fungsi Gauss $f(x) = \eu^{-\pi tx^2}$ bersifat Schwartz, sehingga
\cref{thm:b3:fouriertransform:poisson} berlaku, dan
$\hat f(\xi) = t^{-1/2}\eu^{-\xi^2/4\pi t}$; lalu di $\xi = 2\pi
k$ ruas kanannya menjadi $t^{-1/2}\eu^{-\pi k^2/t}$: yakni
identitas thetanya. Pada $t = 1$:
\[
\sum_{n\in\Z}\eu^{-\pi n^2} = 1 + 2\eu^{-\pi} + 2\eu^{-4\pi} +
\cdots \approx 1 + 0.0864278 + 0.0000070 = 1.0864348,
\]
yang teliti sampai $6$ desimal dengan tiga suku (sebab $\eu^{-9\pi}
\approx 5\cdot10^{-13}$). Pada $t = 10^{-2}$: deret
pendefinisinya memerlukan $\eu^{-\pi n^2/100} < 10^{-7}$, yakni $n \gtrsim
23$ --- kira-kira $47$ suku --- sedangkan deret yang tertransformasi adalah
$10\sum_k\eu^{-100\pi k^2}$, yang sudah pada suku $k = 1$ bernilai
$\sim 10^{-136}$: sehingga satu suku sudah cukup.
\end{solution}

\begin{solution}{exo:b3:fouriertransform:10}
Berlaku $\mathcal Ff \in L^2(\intcc{-\pi}\pi) \subseteq
L^1(\intcc{-\pi}\pi)$ (sebab berukuran berhingga), sehingga pembalikannya memberikan
wakil \omterm{def:b3:topology:continuity}{kontinu} $f(x) =
\frac1{2\pi}\int_{-\pi}^\pi\mathcal Ff(\xi)\eu^{\iu
x\xi}\dd\xi$, dengan
\[
f(n) = \frac1{2\pi}\int_{-\pi}^{\pi}\mathcal
Ff(\xi)\,\eu^{\iu n\xi}\dd\xi = \langle e_{-n}, \mathcal
Ff\rangle
\]
dalam notasi \cref{thm:b3:hilbert:fourier}. Dengan menguraikannya pada
\omterm{def:b3:hilbert:onb}{basis Hilbert} itu: $\mathcal Ff = \sum_nf(n)\,\eu^{-\iu
n\xi}$ di $L^2(\intcc{-\pi}\pi)$. Lalu terapkan operator
kontinu-$L^2$ $\mathcal F^{-1}$ suku demi suku:
\[
\mathcal F^{-1}\bigl(\mathbf
1_{\intcc{-\pi}\pi}\eu^{-\iu n\xi}\bigr)(x)
= \frac1{2\pi}\int_{-\pi}^{\pi}\eu^{\iu\xi(x - n)}\dd\xi
= \frac{\sin\bigl(\pi(x-n)\bigr)}{\pi(x - n)} ,
\]
yang memberikan $f = \sum_nf(n)\operatorname{sinc}(\cdot - n)$ di
$L^2$: sehingga sinyal terbatas pita ditentukan oleh cuplikan
bilangan bulatnya --- yakni teorema pencuplikan Shannon.
\end{solution}

\begin{solution}{exo:b3:fouriertransform:11}
(a) Lewat pembalikan pada $\mathcal S$: $f(x) =
\frac1{2\pi}\int\hat f(\xi)\eu^{\iu x\xi}\dd\xi =
\frac1{2\pi}(\mathcal F\hat f)(-x)$, yakni $(\mathcal
F^2f)(x) = 2\pi f(-x)$. Dengan menerapkannya dua kali: $\mathcal F^4f =
2\pi\,\mathcal F^2f(-\cdot) = (2\pi)^2f$.

(b) Jika $\mathcal Ff = \lambda f$ dengan $f \neq 0$: maka
$(2\pi)^2f = \mathcal F^4f = \lambda^4f$, sehingga $\lambda^4 =
(2\pi)^2$: jadi $\lambda \in \{\pm\sqrt{2\pi},
\pm\iu\sqrt{2\pi}\}$. Sedangkan fungsi Gauss $g(x) = \eu^{-x^2/2}$ mempunyai
$\hat g = \sqrt{2\pi}\,g$
(\cref{ex:b3:fouriertransform:gaussian} di $a = \frac12$):
yakni fungsi eigen untuk $+\sqrt{2\pi}$.

(c) Berlaku $\mathcal F^2f = 2\pi f(-\cdot)$ yang sama dengan $\pm2\pi f$
menurut paritasnya. Untuk $h(x) = x\eu^{-x^2/2}$:
dengan menurunkan $\hat g(\xi) = \sqrt{2\pi}\eu^{-\xi^2/2}$
memakai aturan $\widehat{xf} = \iu\frac{\dd}{\dd\xi}\hat f$:
\[
\hat h(\xi) = \iu\,\frac{\dd}{\dd\xi}\bigl(\sqrt{2\pi}
\eu^{-\xi^2/2}\bigr) = -\iu\sqrt{2\pi}\,\xi\eu^{-\xi^2/2}
= -\iu\sqrt{2\pi}\,h(\xi) :
\]
yakni sebuah fungsi eigen untuk $-\iu\sqrt{2\pi}$. (Fungsi Hermitenya
melanjutkan polanya, dengan berputar melalui keempat
nilai eigennya --- yakni jam Fourier diskretnya.)

\end{solution}

\begin{solution}{exo:b3:fouriertransform:12}
(a) Berlaku $\tilde f \in L^2$ dengan $\norm{\tilde f}_2 = \norm f_2$;
sehingga \cref{exo:b3:lp:6} (dengan eksponen sekawan $p = q = 2$) membuat
$A_f = f * \tilde f$ terbatas dan \omterm{def:b3:topology:continuity}{kontinu} seragam, dengan
\[
A_f(x) = \int f(y)\,\overline{f(y - x)}\,\dd y,
\qquad A_f(0) = \norm f_2^2,
\qquad \abs{A_f(x)} \leq \norm f_2\,\norm{f(\cdot -
x)}_2 = A_f(0)
\]
lewat Cauchy--Schwarz.

(b) Untuk $f \in L^1\cap L^2$: $\tilde f \in L^1$ pula, dan
teorema konvolusinya memberikan $\widehat{A_f} = \hat f\,
\widehat{\tilde f}$; lalu dengan menghitungnya, $\widehat{\tilde f}(\xi) =
\int\overline{f(-x)}\eu^{-\iu\xi x}\dd x =
\overline{\int f(u)\eu^{-\iu\xi u}\dd u} =
\overline{\hat f(\xi)}$: sehingga $\widehat{A_f} = \abs{\hat f}^2
\geq 0$.

(c) Bila $\abs{\hat f}^2 \in L^1$, pembalikannya berlaku bagi
$A_f$ yang \omterm{def:b3:topology:continuity}{kontinu} (sebab transformasinya \omterm{def:b3:lebesgue:l1}{terintegralkan};
\cref{thm:b3:fouriertransform:inversion}):
\[
A_f(x) = \frac1{2\pi}\int\abs{\hat f(\xi)}^2
\eu^{\iu x\xi}\,\dd\xi .
\]
Untuk $f = \mathbf 1_{\intcc{-1/2}{1/2}}$: $\hat f(\xi) =
\frac{2\sin(\xi/2)}\xi = \frac{\sin(\xi/2)}{\xi/2}$, dan di
$x = 0$:
\[
1 = \norm f_2^2 = \frac1{2\pi}\int_\R
\Bigl(\frac{\sin(\xi/2)}{\xi/2}\Bigr)^2\dd\xi
= \frac1{2\pi}\cdot2\int_\R\Bigl(\frac{\sin u}u\Bigr)^2\dd u
\]
(dengan $\xi = 2u$), yakni $\int_\R\bigl(\frac{\sin u}u\bigr)^2\dd
u = \pi$ --- yakni integral kesayangan Plancherel, yang dipulihkan lewat
autokorelasinya.
\end{solution}

\begin{solution}{pb:b3:fouriertransform:1}
\textbf{1.} Dengan mentransformasikan persamaannya terhadap $x$ (secara formal):
$\partial_t\hat u(t,\xi) = \widehat{\partial^2_{xx}u} =
(\iu\xi)^2\hat u = -\xi^2\hat u$, yakni persamaan diferensial biasa terhadap $t$ untuk setiap
frekuensinya: $\hat u(t,\xi) = \eu^{-t\xi^2}\hat f(\xi)$. Karena
$\eu^{-t\xi^2} = \hat g_t$
(\cref{exo:b3:fouriertransform:4}), hasil kalinya adalah
transformasi $g_t * f$.

\textbf{2.} Berlaku $\abs{u(t,x)} \leq \norm f_\infty\int g_t = \norm
f_\infty$: jadi terdefinisi dengan baik. Pada $[t_0, T]\times[-A, A]$: setiap
turunan campurannya $\partial^m_t\partial^n_xg_t(x - y)$ berupa
polinomial dalam $(x - y)$ dan $t^{-1}$ dikali
$\eu^{-(x-y)^2/4t}$, yang untuk $\abs y \geq 2A$ terbatas oleh
$C\,(1 + y^2)^N\eu^{-(\abs y - A)^2/4T}$, yakni pendominasi \omterm{def:b3:lebesgue:l1}{terintegralkan}
yang tak bergantung pada $(t, x)$ di jendela itu (dan terbatas
untuk $\abs y \leq 2A$): sehingga pendiferensialan berulang di bawah
integralnya (\cref{thm:b3:lebesgue:paramdiff}) berlaku: jadi $u \in
\mathcal C^\infty(\intoo0\infty\times\R)$.

\textbf{3.} Dengan $g_t(x) = \frac1{2\sqrt{\pi
t}}\eu^{-x^2/4t}$:
\[
\partial_tg_t = g_t\Bigl(\frac{x^2}{4t^2} - \frac1{2t}\Bigr)
= \partial^2_{xx}g_t
\]
(turunkanlah dua kali terhadap $x$: $\partial_xg_t = -\frac
x{2t}g_t$, sehingga $\partial^2_{xx}g_t = \bigl(\frac{x^2}{4t^2} -
\frac1{2t}\bigr)g_t$). Lalu menurut pertanyaan 2 turunannya lewat
di bawah integralnya: sehingga $\partial_tu = \partial^2_{xx}u$.

\textbf{4.} Berlaku $u(t,x) - f(x) = \int g_t(y)\bigl(f(x - y) -
f(x)\bigr)\dd y$. Diberikan sebuah \omterm{def:b3:topology:compact}{kompak} $K$ dan $\varepsilon$:
\omterm{def:b3:topology:continuity}{kekontinuan} seragam $f$ pada sebuah \omterm{def:b3:topology:topology}{persekitaran} $K$ memberikan
$\delta$ dengan $\abs{f(x-y) - f(x)} < \varepsilon$ untuk $x \in
K$ dan $\abs y \leq \delta$; sedangkan ekornya menyumbang $\leq 2\norm
f_\infty\int_{\abs y > \delta}g_t(y)\dd y = 2\norm
f_\infty\,\P$ yakni massa di luar $\delta$, yang bernilai
$\frac2{\sqrt\pi}\int_{\delta/2\sqrt t}^\infty\eu^{-z^2}\dd z
\to 0$ saat $t \to 0$. Untuk $f \in L^p$: $\norm{u(t) - f}_p \leq
\int g_t(y)\norm{\tau_yf - f}_p\dd y$ (lewat Minkowski/Jensen seperti pada
\cref{thm:b3:lp:regularization}), lalu pisahkan dengan cara yang sama memakai
\cref{thm:b3:lp:density}(3).

\textbf{5.} Pemulusan seketikanya adalah pertanyaan 2 (sebab $u(t)$ bersifat
$\mathcal C^\infty$ untuk $t > 0$ tanpa memakai kemulusan $f$
apa pun). Untuk $f = \mathbf 1_{\intoo0\infty}$:
\[
u(t, x) = \int_0^\infty g_t(x - y)\dd y
= \frac1{\sqrt\pi}\int_{-x/2\sqrt t}^{\infty}\eu^{-z^2}\dd z
= \frac12\Bigl(1 +
\operatorname{erf}\Bigl(\frac{x}{2\sqrt t}\Bigr)\Bigr),
\qquad \operatorname{erf}(s) =
\frac2{\sqrt\pi}\int_0^s\eu^{-z^2}\dd z :
\]
yakni sebuah tangga yang termuluskan dengan zona peralihan yang melebar seperti $\sqrt t$
(dengan profil pada $t_1 < t_2 < t_3$: yakni landaian yang makin datar melalui
$(0, \frac12)$).

\textbf{6.} Integrannya $g_t(x-y)f(y)$ bernilai $\geq 0$ dan
kernelnya positif sejati: sehingga $u(t,x) = 0$ akan memaksa $f = 0$
hampir di mana-mana. Sedangkan kemonotonan terhadap $f$ adalah kemonotonan integralnya. Sebuah
titik yang $f = 0$-nya pada sebuah selang tetap mempunyai $u(t, \cdot) > 0$
di sana untuk setiap $t > 0$: sehingga panas merambat dengan laju tak berhingga
(sebab kepositifan di mana pun terasa di mana-mana seketika).

\textbf{7.} Lewat Tonelli (sebab $g_t(x-y)\abs{f(y)}$ \omterm{def:b3:lebesgue:l1}{terintegralkan} pada
$\R^2$): $\int u(t,x)\dd x = \int f(y)\bigl(\int g_t(x -
y)\dd x\bigr)\dd y = \int f$.

\textbf{8.} Lewat Plancherel: $2\pi\norm{u(t)}_2^2 =
\int\eu^{-2t\xi^2}\abs{\hat f(\xi)}^2\dd\xi$, yang tak naik terhadap
$t$ (secara titik demi titik), dan sejati kecuali $\hat f = 0$ hampir di mana-mana (yakni $= f =
0$), dengan limit $0$ saat $t\to\infty$ lewat kekonvergenan terdominasi. Dan
$\norm{u(t)}_\infty \leq \norm{g_t}_\infty\norm f_1 =
\frac{\norm f_1}{2\sqrt{\pi t}} \to 0$.

\textbf{9.} Untuk $\varphi \in \mathcal C_c^\infty$, $t \mapsto
\langle\varphi, \hat u(t)\rangle$ bersifat $\mathcal C^1$ dengan
turunan $\langle\varphi, \partial_t\hat u\rangle =
\langle\varphi, \widehat{\partial_{xx}u}\rangle =
\langle\xi^2\varphi\dots\rangle$ --- lebih tepatnya,
$\widehat{\partial^2_{xx}u} = -\xi^2\hat u$ memindahkan
persamaannya. Lalu untuk hampir setiap $\xi$, fungsi $t \mapsto \eu^{t\xi^2}\hat u(t,\xi)$
yang \omterm{def:b3:topology:continuity}{kontinu} mutlak berturunan
$\eu^{t\xi^2}(\xi^2\hat u + \partial_t\hat u) = 0$ dalam
arti \omterm{def:b3:lebesgue:l1}{terintegralkannya}: sehingga ia konstan, dan dengan membiarkan $t \to 0$
(sebab $\hat u(t) \to \hat f$ di $L^2$, jadi hampir di mana-mana sepanjang sebuah subbarisan):
$\hat u(t, \xi) = \eu^{-t\xi^2}\hat f(\xi)$ hampir di mana-mana. Jadi dua
solusi pada kelas itu bertransformasi sama: sehingga keduanya
sama.

\textbf{10.} Bahwa $f = g_s * h$ dengan $h \in L^2$ memaksa $\hat f =
\eu^{-s\xi^2}\hat h$, yakni $\hat h = \eu^{s\xi^2}\hat f \in
L^2$. Ambillah $\hat f(\xi) = \eu^{-\abs\xi}$: maka $f(x) =
\frac1\pi\cdot\frac1{1 + x^2}$
(\cref{exo:b3:fouriertransform:1}, lewat pembalikan), yakni fungsi $L^2$
yang mulus \omterm{prop:b3:galois:perfect}{sempurna}; padahal $\eu^{2s\xi^2 - 2\abs\xi} \to
\infty$: sehingga $\eu^{s\xi^2}\hat f \notin L^2$ untuk setiap $s > 0$.
Jadi profil Cauchynya \emph{tak pernah} merupakan hasil difusi
sebelumnya.

\textbf{11.} Semigrup panasnya mengalikan transformasinya dengan
$\eu^{-t\xi^2}$, yang tak lenyap di mana pun: jadi injektif ---
sehingga secara formal tak ada keterangan yang dimusnahkan. Tetapi jangkauannya terdiri atas
fungsi yang transformasinya meluruh seperti $\eu^{-t\xi^2}$: yakni
subruang $L^2$ yang mungil, padat, tetapi sejati (sebab pertanyaan 10 menunjukkan
bahwa fungsi yang sangat baik pun terletak di luarnya). Membalikkannya akan
menguatkan frekuensi $\xi$ dengan $\eu^{t\xi^2}$: yang tak terbatas,
sehingga tak stabil terhadap gangguan apa pun. Jadi difusinya
tak terbalikkan bukan karena pemetaannya lupa, melainkan karena
balikannya tak dapat \omterm{def:b3:topology:continuity}{kontinu} --- yakni anak panah waktu yang terbuat dari
analisis fungsional.

\textbf{12.} Berlaku $\hat f \in L^2(\intcc{-\Omega}\Omega)
\subseteq L^1$ (lewat Cauchy--Schwarz pada selang terbatas), sehingga
$F(x) = \frac1{2\pi}\int_{-\Omega}^\Omega\hat
f(\xi)\eu^{\iu x\xi}\dd\xi$ terdefinisi di mana-mana, dan
pendiferensialan di bawah integralnya (yang terdominasi oleh
$\Omega^k\abs{\hat f} \in L^1$ pada pitanya) membuatnya
$\mathcal C^\infty$ dengan $\abs{F^{(k)}} \leq
\frac{\Omega^k}{2\pi}\norm{\hat f}_{L^1}$ di mana-mana. Dan $F
= f$ hampir di mana-mana: sebab kedua ruasnya bertransformasi sama, dan
transformasinya injektif pada $L^2$
(\cref{thm:b3:fouriertransform:plancherel} beserta perluasan
$L^2$-nya).

\textbf{13.} Eksponensial $\xi \mapsto \eu^{-\iu
n\pi\xi/\Omega}$ dengan $n \in \Z$ membentuk \omterm{def:b3:hilbert:onb}{basis Hilbert}
$L^2(\intcc{-\Omega}\Omega)$
(\cref{thm:b3:hilbert:fourier}, yang diskalakan ulang). Koefisien
$\hat f$ sepanjang yang ke-$n$ adalah
\[
\frac1{2\Omega}\int_{-\Omega}^\Omega\hat f(\xi)\,
\eu^{\iu n\pi\xi/\Omega}\dd\xi
= \frac{2\pi}{2\Omega}\cdot
\frac1{2\pi}\int_{-\Omega}^{\Omega}\hat f(\xi)\,
\eu^{\iu(n\pi/\Omega)\xi}\dd\xi
= \frac\pi\Omega\,f\Bigl(\frac{n\pi}\Omega\Bigr),
\]
menurut rumus pertanyaan 12 di $x = \frac{n\pi}\Omega$: sehingga
uraian yang dinyatakan itu berlaku di $L^2$ pitanya.

\textbf{14.} Sisipkan uraiannya ke dalam rumus pembalikan
pertanyaan 12; sedangkan penukaran jumlah dan integralnya adalah
\omterm{def:b3:topology:continuity}{kekontinuan} pemasangan $L^2$ terhadap $\frac1{2\pi}
\eu^{\iu x\xi}\mathbf 1_{\abs\xi\leq\Omega}$ (yang bernorma-$L^2$
$\frac{\sqrt{2\Omega}}{2\pi}$ dan tak bergantung pada $x$ ---
sehingga keseragamannya):
\[
f(x) = \sum_nf\Bigl(\frac{n\pi}\Omega\Bigr)\cdot
\frac1{2\Omega}\int_{-\Omega}^\Omega
\eu^{\iu(x - n\pi/\Omega)\xi}\dd\xi
= \sum_nf\Bigl(\frac{n\pi}\Omega\Bigr)
\operatorname{sinc}(\Omega x - n\pi),
\]
sebab $\frac1{2\Omega}\int_{-\Omega}^\Omega\eu^{\iu
u\xi}\dd\xi = \frac{\sin(\Omega u)}{\Omega u}$.

\textbf{15.} Dengan membaca perhitungan pertanyaan 14 secara mundur,
transformasi $s_n = \operatorname{sinc}(\Omega\cdot -
n\pi)$ adalah $\hat s_n = \frac\pi\Omega\,\eu^{-\iu
n\pi\xi/\Omega}\,\mathbf 1_{\intcc{-\Omega}\Omega}$.
Lalu Plancherel:
\[
\langle s_n, s_m\rangle = \frac1{2\pi}
\Bigl(\frac\pi\Omega\Bigr)^2\int_{-\Omega}^\Omega
\eu^{\iu(n-m)\pi\xi/\Omega}\dd\xi =
\frac\pi\Omega\,\delta_{nm} :
\]
yakni keluarga ortogonal yang bernorma tetap
$\sqrt{\pi/\Omega}$. Lalu dengan mengambil norma pada uraian
pertanyaan 14: $\norm f_2^2 =
\frac\pi\Omega\sum_n\abs{f(n\pi/\Omega)}^2$.

\textbf{16.} Fungsi $g(x) = \sin(\Omega x)\operatorname{sinc}
(\Omega x) = \frac{\sin^2(\Omega x)}{\Omega x}$ lenyap di
setiap titik kisi $\frac{n\pi}\Omega$ (termasuk $0$, lewat
limitnya) dan tak identik nol. Untuk pitanya: tulislah $g =
\frac1{2\iu}\bigl(\eu^{\iu\Omega x} - \eu^{-\iu\Omega
x}\bigr)\operatorname{sinc}(\Omega x)$; sebab modulasi oleh
$\eu^{\pm\iu\Omega x}$ menggeser transformasinya sejauh
$\mp\Omega$, sehingga $\hat g$ berpendukung di
$\intcc{-2\Omega}{2\Omega}$ (sungguh di gabungan dua
pita yang tergeser): jadi $g \in PW_{2\Omega}$, yang tak terlihat oleh
pencuplikan berlaju-$\Omega$ --- yakni pelipatan spektrum yang menjelma.

\textbf{17.} Menurut pertanyaan 15 cuplikannya membawa energinya
secara demokratis: $\norm f^2 = \frac\pi\Omega\sum
\abs{f(n\pi/\Omega)}^2$. Jika energi sinyalnya di luar
jendela waktu $\intcc{-T/2}{T/2}$ bernilai $\leq \varepsilon^2\norm
f^2$, maka cuplikan di luar jendelanya memenuhi (sampai suku
batas yang dikendalikan batas seragam pertanyaan 12)
$\frac\pi\Omega\sum_{\abs{n\pi/\Omega} > T/2}
\abs{f(n\pi/\Omega)}^2 \approx \norm{f\,\mathbf 1_{\abs x >
T/2}}^2 \leq \varepsilon^2\norm f^2$: sehingga memancung
deret pencuplikannya menjadi $\approx \frac{\Omega T}\pi$
indeks di dalam jendelanya merekonstruksi $f$ sampai galat nisbi
$\approx\varepsilon$. Karena itu hasil kali waktu dan lebar pita
$\frac{\Omega T}{\pi}$ mencacah derajat kebebasan real
efektif sinyalnya --- yakni kaidah di balik setiap
format audio.

\textbf{18.} (a) Fungsi $\operatorname{sinc}(\Omega x)$ bercuplikan
$f(n\pi/\Omega) = \operatorname{sinc}(n\pi) =
\delta_{n0}$: sehingga deretnya menyusut ke suku $n = 0$-nya, yakni
$\operatorname{sinc}(\Omega x)$ --- jadi teoremanya mereproduksi
kernelnya sendiri. (b) Jika $\hat f$ berpendukung di
$\intcc{-\Omega'}{\Omega'} \subseteq
\intcc{-\Omega}\Omega$, maka setiap langkah pertanyaan 13--14 berjalan
kata demi kata dengan pita $\Omega$ yang lebih besar (sebab uraian
$\hat f$ pada selang yang lebih besar tetap sah):
sehingga mencuplik lebih cepat daripada laju Nyquistnya sendiri tak mengubah apa pun
pada rekonstruksinya --- jadi pencuplikan berlebih tak berbahaya, dan pada
praktiknya bermanfaat (sebab kernel rekonstruksi yang meluruh lebih cepat
lalu dapat dipakai).

\textbf{19.} Uraikan $\eu^{-\iu\xi x} =
\sum_k\frac{(-\iu\xi x)^k}{k!}$ di dalam integralnya; lalu pada
$\intcc{-A}A$ deretnya konvergen secara normal
(sebab $\sum_k\frac{\abs{\xi}^kA^k}{k!}\abs f \in L^1$), sehingga
pengintegralan suku demi suku bersifat sah:
\[
\hat f(\xi) = \sum_{k\geq0}\frac{(-\iu\xi)^k}{k!}m_k,
\qquad \abs{m_k} \leq A^k\norm f_1 .
\]
Batas itu membuat deretnya konvergen untuk setiap $\xi$ yang kompleks;
lalu di sekitar sembarang titik $\xi_0$, pengelompokan ulangnya (lewat kekonvergenan mutlaknya)
memberikan deret pangkat dalam $\xi - \xi_0$: sehingga $\hat f$ bersifat
analitik real dengan jari-jari tak berhingga di mana-mana.

\textbf{20.} Misalkan $g$ analitik real pada $\R$ (yakni deret Taylornya
konvergen ke $g$ di dekat setiap titik) dan $Z = \{\xi :
g^{(k)}(\xi) = 0\ \forall k\}$. Maka $Z$ tertutup (sebab irisan
himpunan tertutup); dan ia terbuka, sebab di $\xi_0 \in Z$ uraian
Taylor lokal $g$ adalah deret nol, sehingga $g$ lenyap
secara identik di dekat $\xi_0$, beserta semua turunannya. Jika
$g$ lenyap pada sebuah selang, maka $Z \neq \varnothing$; lalu menurut
keterhubungan $\R$, $Z = \R$: sehingga $g \equiv 0$. Kini jika $f
\neq 0$ berpendukung \omterm{def:b3:topology:compact}{kompak} bersama $\hat f$: maka pertanyaan
19 membuat $\hat f$ analitik real, yang lenyap di luar sebuah \omterm{def:b3:topology:compact}{kompak},
jadi pada selang: sehingga $\hat f \equiv 0$, jadi $f = 0$ hampir di mana-mana lewat
keinjektifannya --- yang bertentangan. Demikian pula $f \in
PW_\Omega$ yang tak nol tak dapat berpendukung \omterm{def:b3:topology:compact}{kompak} (dengan menukar peran
$f$ dan $\hat f$ lewat pembalikannya): sehingga sinyal terbatas pita tak pernah
mati, dan sinyal terbatas waktu menempati spektrum yang tak terbatas.

\textbf{21.} Untuk $f = \eu^{-ax^2}$: $\norm f_2^2 =
\sqrt{\frac\pi{2a}}$ dan $\int x^2\abs f^2 =
\frac1{4a}\sqrt{\frac{\pi}{2a}}$ (lewat momen kedua Gaussnya);
sedangkan $\hat f = \sqrt{\frac\pi a}\,\eu^{-\xi^2/4a}$
(\cref{ex:b3:fouriertransform:gaussian}) dan
\[
\frac1{2\pi}\int\xi^2\abs{\hat f}^2\dd\xi =
\frac1{2\pi}\cdot\frac\pi a\int\xi^2
\eu^{-\xi^2/2a}\dd\xi = \frac1{2a}\cdot a\sqrt{2\pi a}
= \frac{\sqrt{2\pi a}}2 .
\]
Hasil kali ternormalkannya: $\frac1{4a}\sqrt{\frac\pi{2a}}\cdot
\frac{\sqrt{2\pi a}}2\big/\frac{\pi}{2a} = \frac14$,
yang tak bergantung pada $a$. Penskalaannya menerangkan ketetapannya: sebab mengganti
$f$ dengan $f(\lambda\cdot)$ mengalikan $\int x^2\abs f^2/\norm
f^2$ dengan $\lambda^{-2}$ dan $\frac1{2\pi}\int\xi^2\abs{\hat
f}^2/\norm f^2$ dengan $\lambda^{2}$: sehingga hasil kalinya invarian terhadap
pendilatan, dan fungsi Gaussnya membentuk satu \omterm{def:b3:groups:action}{orbit} pendilatan.

\textbf{22.} Dengan $\abs f^2$ sebagai \omterm{ex:b3:lebesgue:gamma}{kepadatan} posisi dan
$\frac1{2\pi}\abs{\hat f}^2$ \omterm{ex:b3:lebesgue:gamma}{kepadatan} momentum sebuah
keadaan kuantum (sedangkan satuan fisikanya menyisipkan $\hbar$),
\cref{exo:b3:fouriertransform:8} berbunyi $\sigma_x\sigma_p
\geq \frac\hbar2$: sehingga tak ada keadaan yang tajam pada kedua teramatnya.
Di sepanjang bab ini, satu hukum mengenakan lima jubah: panas seketika
memuluskan sebab $\eu^{-t\xi^2}$ memusnahkan frekuensi tinggi
(Bagian II); alirannya tak dapat berjalan mundur sebab memulihkannya
bersifat tak terbatas (Bagian IV); sinyal terbatas pita cukup kaku
untuk hidup pada kisi terbilang (Bagian V); tak ada fungsi yang
mengalahkan lantai Heisenberg; dan tak ada fungsi yang berpendukung \omterm{def:b3:topology:compact}{kompak}
pada kedua sisi transformasinya (pertanyaan 19--20).
Jadi apa yang dilakukan $\hat f$ di tak hingga mengatur apa yang boleh dilakukan $f$
di mana pun.

\textbf{23.} Baik $g_t$ maupun $g_s$ berada di $L^1$ dengan
$\widehat{g_t}(\xi) = \eu^{-t\xi^2}$ (lewat perhitungan pertanyaan
1), sehingga teorema konvolusinya memberikan
$\widehat{g_t * g_s} = \eu^{-t\xi^2}\eu^{-s\xi^2} =
\eu^{-(t+s)\xi^2} = \widehat{g_{t+s}}$; dan dua fungsi $L^1$
yang bertransformasi sama bersesuaian hampir di mana-mana (lewat keinjektifannya, yaitu lewat
teorema pembalikannya --- dan di sini kedua ruasnya \omterm{def:b3:topology:continuity}{kontinu}, sehingga
keduanya bersesuaian di mana-mana): jadi $g_t * g_s = g_{t+s}$. Karena itu
$u(t + s) = g_{t+s} * f = g_s * (g_t * f) = g_s * u(t)$
(lewat keasosiatifan konvolusinya dan Tonelli). Untuk kelog-cembungannya: misalkan
$N(t) = \norm{u(t)}_2^2 = \frac1{2\pi}\int
\eu^{-2t\xi^2}\abs{\hat f}^2\dd\xi$ (lewat Plancherel, pertanyaan
8). Untuk $t = \frac{t_1 + t_2}2$, tulislah
\[
\eu^{-2t\xi^2}\abs{\hat f}^2
= \Bigl(\eu^{-2t_1\xi^2}\abs{\hat f}^2\Bigr)^{1/2}
\Bigl(\eu^{-2t_2\xi^2}\abs{\hat f}^2\Bigr)^{1/2},
\]
lalu Cauchy--Schwarz memberikan $N\bigl(\frac{t_1+t_2}2\bigr) \leq
\sqrt{N(t_1)\,N(t_2)}$: sehingga $\ln N$ cembung di titik tengahnya, dan
karena \omterm{def:b3:topology:continuity}{kontinu} (lewat kekonvergenan terdominasi terhadap $t$), ia cembung; jadi
demikian pula $\ln\norm{u(t)}_2 = \frac12\ln N(t)$. Yakni peluruhan berlogaritma
cembung: sehingga aliran panasnya tak dapat kehilangan energi dalam satu ledakan lalu
mandek.

\textbf{24.} Momen kernelnya: $\int g_t = 1$ (lewat pertanyaan
7 dengan $f = g_s$, atau langsung lewat integral Gaussnya), $\int
x\,g_t(x)\dd x = 0$ (sebab integrannya gasal), dan, dengan menyubstitusikan $x =
2\sqrt t\,v$,
\[
\int_\R x^2g_t(x)\,\dd x
= \frac{4t}{\sqrt\pi}\int_\R v^2\eu^{-v^2}\dd v = 2t .
\]
Lalu dengan menyubstitusikan $x = z + y$ pada konvolusinya dan mengamati bahwa
$\iint(\abs z + \abs y)^2g_t(z)f(y)\,\dd z\,\dd y < \infty$
(sebab masing-masing $\int\abs z^kg_t$ dan $\int\abs y^kf$ berhingga untuk $k
\leq 2$, dengan memakai $\abs y \leq \frac{1 + y^2}2$), Fubini dan
Tonelli berlaku bagi integral momen di bawah ini:
\[
\int x\,u(t,x)\dd x = \iint (z + y)\,g_t(z)f(y)\,\dd z\,\dd
y = 0\cdot\!\int\! f + 1\cdot\!\int\! yf(y)\dd y,
\]
yakni klaim pertamanya; dan
\[
\iint (z+y)^2g_t(z)f(y)\,\dd z\,\dd y
= 2t\int f + 2\cdot0\cdot\!\int\! yf + \int y^2f(y)\dd y ,
\]
yakni yang kedua. Jadi meannya menjumlah, ragamnya menjumlah, dan kernelnya
menyumbang mean $0$ serta ragam $2t$: sehingga setelah waktu $t$
panasnya telah menyebar selebar orde $\sqrt{2t}$ ---
yakni jarak tumbuh seperti akar kuadrat waktunya, yaitu tanda tangan
difusinya (dan \omterm{def:b3:topology:pathconnected}{lintasan} Brown pada
\cref{ch:b3:probability}).

\textbf{25.} Pada sisi transformasinya: $\hat f(\xi) =
\sqrt\pi\,\eu^{-\xi^2/4}$, sehingga $\hat u(t,\xi) =
\sqrt\pi\,\eu^{-(t + \frac14)\xi^2}$, yang merupakan transformasi
$(1 + 4t)^{-1/2}\exp\bigl(-x^2/(1+4t)\bigr)$ (lewat
kamus Gaussnya $\eu^{-ax^2} \mapsto
\sqrt{\pi/a}\,\eu^{-\xi^2/4a}$ dengan $a = \frac1{1+4t}$):
yakni bentuk tertutupnya. Untuk pemeriksaan langsungnya, dengan $\sigma = 1 + 4t$:
\[
\partial_tu = \sigma^{-1/2}\eu^{-x^2/\sigma}
\Bigl(-\frac2\sigma + \frac{4x^2}{\sigma^2}\Bigr)
= \partial^2_{xx}u ,
\]
dengan kedua ruasnya dihitung dari $\partial_xu =
-\frac{2x}\sigma\,u$. Untuk kekekalannya: $\int u(t) =
\sigma^{-1/2}\sqrt{\pi\sigma} = \sqrt\pi$ untuk setiap $t$.
Untuk pelesapannya:
\[
\norm{u(t)}_2^2 = \frac1\sigma\int\eu^{-2x^2/\sigma}\dd x
= \frac1\sigma\sqrt{\frac{\pi\sigma}2}
= \sqrt{\frac\pi2}\,(1+4t)^{-1/2},
\]
sehingga $\norm{u(t)}_2 = (\pi/2)^{1/4}(1+4t)^{-1/4}$,
yang tak naik dengan logaritma yang cembung (pertanyaan 23); dan
norma $L^2$-nya meluruh seperti $t^{-1/4}$, tepat separuh
eksponen $t^{-1/2}$ milik $\norm{u(t)}_\infty$ --- yang konsisten
dengan $\norm u_2^2 \leq \norm u_\infty\norm u_1$ dan
kekekalan $\norm u_1$. Untuk ragamnya: $\int x^2u(t) =
\frac1\sigma\cdot\frac{\sigma^{3/2}\sqrt\pi}2 =
\frac{\sqrt\pi}2\,(1 + 4t) = \int x^2f + 2t\sqrt\pi$, seperti yang
diramalkan pertanyaan 24 (sebab $\int x^2f = \frac{\sqrt\pi}2$ dan $\int
f = \sqrt\pi$). Pada $t = 6$: $\sigma = 25$, dengan tinggi puncak
$u(6, 0) = \frac15$ terhadap $u(0,0) = 1$, skala lebarnya
$\sqrt\sigma = 5$ kali skala awalnya, dan $\int u =
\sqrt\pi \approx 1.7725$ sepanjang waktu: sehingga bercaknya lima kali
lebih rendah, lima kali lebih lebar, dan tak sekalori pun hilang.

\end{solution}
